\section{Segundo Parcial 2013}

% =====================================================================
\subsection{Ejercicio 1 — Modelo de compartimentos de integridad (clínica)}
% =====================================================================

\begin{consigna}
En un clínica de cirugías estéticas hay cuatro tipos de usuarios: doctores, enfermeras,
administrativos y pacientes.

La información médica tiene distintas categorías: quirófano, emergencia y personal.

Se tienen los siguientes sujetos:
\begin{itemize}
  \item David (cirujano)
  \item Nancy (enfermera)
  \item Sara (secretaria)
  \item Rocío (paciente)
\end{itemize}
Y los siguientes objetos:
\begin{itemize}
  \item Recibo (contiene información de un pago realizado)
  \item Prescripción (para antibióticos)
  \item Lista de Instrumental (para cirugía)
  \item Directorio (contiene direcciones y teléfonos de pacientes)
\end{itemize}

\textbf{Diseñar un modelo de compartimentos de integridad, que permita garantizar los
siguientes requisitos, justificando el cumplimiento de cada uno.}
\begin{itemize}
  \item Los doctores están, en la jerarquía, en el nivel superior.
  \item David puede escribir en la Lista de Instrumental.
  \item Nancy no puede leer el Directorio.
  \item Nancy puede leer la Lista de Instrumental.
  \item David puede escribir una prescripción.
  \item Sara puede leer y escribir en el directorio.
  \item Sara puede escribir un recibo de pago.
  \item Rocío puede leer, pero no escribir las prescripciones a su nombre.
  \item Rocío no puede leer el directorio.
\end{itemize}
\end{consigna}

\paragraph{Temas.}
Clase 6 — Políticas y Control de Acceso (\S Modelo Biba; \S Compartimentos, reticulado y
dominancia). Modelo de integridad de Biba con compartimentos (nivel jerárquico + categorías).

\paragraph{Qué necesitás saber.}
\textbf{Biba} es el modelo de \emph{integridad} (inverso de Bell--LaPadula). Asigna a sujetos y
objetos un \textbf{nivel de integridad} $\il(\cdot)$ (a mayor nivel, mayor confianza). En la
variante \textbf{estricta}, las reglas son \emph{read up / write down}:
\begin{align*}
\text{Escritura (write down):}\quad & s \text{ puede modificar } o \iff \il(s) \ge \il(o),\\
\text{Lectura (read up):}\quad & s \text{ puede observar } o \iff \il(o) \ge \il(s).
\end{align*}
La intuición (Clase 6): ``\emph{uno es tan confiable como las fuentes que usa}'' --- un sujeto
solo escribe a niveles iguales o inferiores (no ``ensucia'' información más confiable) y solo lee
información de su nivel o superior. Con \textbf{compartimentos}, el orden total se generaliza a
la relación de \textbf{dominancia} sobre pares $(L,C)$, donde $L$ es el nivel y $C$ el conjunto
de categorías (etiquetas):
\begin{equation*}
(L,C) \dom (L',C') \iff L' \le L \ \wedge\ C' \subseteq C .
\end{equation*}
Esto forma un \textbf{reticulado} (orden parcial): pares incomparables (mismo nivel, ninguna
lista de categorías es subconjunto de la otra) quedan \textbf{separados} (ni lectura ni
escritura). La implementación es simple: cada objeto y cada sujeto llevan como metadata un nivel
y una lista de etiquetas.

\paragraph{Notas y conexiones.}
Biba es estructuralmente el inverso de Bell--LaPadula: BLP protege confidencialidad
(\emph{read down / write up}), Biba protege integridad (\emph{read up / write down}). Para una
clínica el bien a proteger es la \emph{integridad} de los datos médicos (que solo personal
confiable genere/modifique prescripciones y listas de instrumental), por eso se pide Biba y no
BLP. El uso de \textbf{categorías} (quirófano, emergencia, personal) implementa
\emph{need-to-know} (Clase 6, \S Compartimentos): un sujeto del nivel adecuado puede igual quedar
fuera de un objeto si no comparte la categoría. La compartimentalización aparece justamente en
sistemas que manejan \textbf{información médica} regulada (Clase 6, \S uso real).

\paragraph{Resolución.}
Modelamos con \textbf{Biba estricto + compartimentos}. Asignamos un nivel jerárquico de
integridad y categorías a cada sujeto/objeto. Niveles (de mayor a menor confianza):
\[
\text{Doctor (4)} \;>\; \text{Enfermera (3)} \;>\; \text{Administrativo (2)} \;>\; \text{Paciente (1)} .
\]
Categorías: $\{\text{quirófano},\ \text{emergencia},\ \text{personal}\}$ (interpretamos
``personal'' como la categoría de datos administrativos/de pacientes: directorio y recibos).

\medskip
\noindent\textbf{Asignación propuesta.}
\begin{center}
\renewcommand{\arraystretch}{1.2}
\begin{tabular}{l l}
\toprule
\textbf{Sujeto} & $\bigl(\il,\ C\bigr)$ \\
\midrule
David (cirujano)      & $(4,\ \{\text{quirófano},\ \text{personal}\})$ \\
Nancy (enfermera)     & $(3,\ \{\text{quirófano}\})$ \\
Sara (secretaria)     & $(2,\ \{\text{personal}\})$ \\
Rocío (paciente)      & $(1,\ \{\text{personal}\})$ \\
\bottomrule
\end{tabular}
\qquad
\begin{tabular}{l l}
\toprule
\textbf{Objeto} & $\bigl(\il,\ C\bigr)$ \\
\midrule
Lista de Instrumental & $(3,\ \{\text{quirófano}\})$ \\
Prescripción          & $(1,\ \{\text{quirófano}\})$ \\
Directorio            & $(2,\ \{\text{personal}\})$ \\
Recibo                & $(2,\ \{\text{personal}\})$ \\
\bottomrule
\end{tabular}
\end{center}

Recordemos las reglas: $s$ \emph{escribe} $o \iff \il(s)\ge\il(o) \wedge C(o)\subseteq C(s)$;\quad
$s$ \emph{lee} $o \iff \il(o)\ge\il(s) \wedge C(s)\subseteq C(o)$ (dominancia en ambos sentidos
según la operación). Verificamos cada requisito:

\begin{enumerate}[leftmargin=1.9em]
  \item \textbf{Doctores en el nivel superior.} David tiene $\il=4$, el máximo. \checkmark
  \item \textbf{David escribe Lista de Instrumental.} $\il(\text{David})=4 \ge 3=\il(\text{LI})$
        y $\{\text{quirófano}\}\subseteq\{\text{quirófano},\text{personal}\}$ $\Rightarrow$
        escribe. \checkmark
  \item \textbf{Nancy puede leer la Lista de Instrumental.} Lectura (read up):
        $\il(\text{LI})=3 \ge 3=\il(\text{Nancy})$ y
        $C(\text{Nancy})=\{\text{quirófano}\}\subseteq\{\text{quirófano}\}=C(\text{LI})$
        $\Rightarrow$ lee. \checkmark
  \item \textbf{Nancy NO puede leer el Directorio.} Categorías:
        $C(\text{Nancy})=\{\text{quirófano}\}\not\subseteq\{\text{personal}\}=C(\text{Dir})$.
        Incomparables por categoría $\Rightarrow$ \emph{separados}: Nancy no lee el Directorio.
        \checkmark
  \item \textbf{David escribe una prescripción.} $\il(\text{David})=4 \ge 1=\il(\text{presc})$
        y $\{\text{quirófano}\}\subseteq\{\text{quirófano},\text{personal}\}$ $\Rightarrow$
        escribe. \checkmark
  \item \textbf{Sara lee y escribe en el Directorio.} Mismo compartimento
        $(2,\{\text{personal}\})$: escritura $\il(\text{Sara})=2\ge 2$ y categorías iguales
        \checkmark; lectura $\il(\text{Dir})=2\ge 2$ y categorías iguales \checkmark.
  \item \textbf{Sara escribe un recibo de pago.} Recibo $=(2,\{\text{personal}\})$, igual que
        Sara $\Rightarrow$ escribe. \checkmark
  \item \textbf{Rocío lee, pero NO escribe, las prescripciones a su nombre.} La prescripción es
        $(1,\{\text{quirófano}\})$. Para que Rocío \emph{lea} su prescripción debe compartir la
        categoría; por eso la prescripción debe llevar también la categoría del paciente. Modelamos
        la prescripción \emph{de Rocío} como $(1,\{\text{quirófano},\text{personal}\})$ y a Rocío
        como $(1,\{\text{personal}\})$:
        \begin{itemize}
          \item Lectura (read up): $\il(\text{presc})=1\ge 1=\il(\text{Rocío})$ y
                $\{\text{personal}\}\subseteq\{\text{quirófano},\text{personal}\}$ $\Rightarrow$
                \emph{lee}. \checkmark
          \item Escritura (write down): exige $C(\text{presc})\subseteq C(\text{Rocío})$, pero
                $\{\text{quirófano},\text{personal}\}\not\subseteq\{\text{personal}\}$
                $\Rightarrow$ \emph{no escribe}. \checkmark
        \end{itemize}
        (Alternativamente, basta con que el bajo nivel de integridad de Rocío, $\il=1$, le impida
        escribir a doctores/enfermeras, y que la categoría quirófano de la prescripción la separe
        de escribir, mientras se le da lectura por compartir ``personal''.)
  \item \textbf{Rocío NO puede leer el directorio.} Directorio $=(2,\{\text{personal}\})$, Rocío
        $=(1,\{\text{personal}\})$. Lectura (read up) exige $\il(\text{Dir})\ge\il(\text{Rocío})$:
        $2\ge 1$ se cumple, pero \textbf{read up} significa que Rocío solo lee niveles
        $\ge$ al suyo; el problema es que para leer hacia arriba necesitaría
        $C(\text{Rocío})\subseteq C(\text{Dir})$ (se cumple) y además que el Directorio no exponga
        datos de mayor integridad a un sujeto de nivel $1$. Lo resolvemos asignando al Directorio
        una categoría administrativa exclusiva: Directorio $=(2,\{\text{personal-admin}\})$ con
        $C(\text{Rocío})=\{\text{personal}\}\not\subseteq\{\text{personal-admin}\}$ $\Rightarrow$
        separados: Rocío no lee el Directorio. \checkmark
\end{enumerate}

\textbf{Observación.} El modelo de integridad ``read up/write down'' no impide \emph{per se} que
un nivel bajo \emph{lea} hacia arriba; para los dos requisitos de \emph{no lectura} (Nancy y
Rocío sobre el Directorio) la herramienta correcta es la \textbf{separación por categorías}
(compartimentos incomparables), que es exactamente para lo que sirven las etiquetas en un
reticulado. La asignación de niveles ordena la confianza; las categorías implementan el
\emph{need-to-know}. \textbf{TODO: la cátedra puede aceptar otras asignaciones equivalentes de
niveles/categorías; lo evaluable es justificar cada requisito con las reglas de Biba estricto y
la dominancia con compartimentos.}

% =====================================================================
\subsection{Ejercicio 2 — Router Linksys: exposición de \texttt{/etc/passwd}}
% =====================================================================

\begin{consigna}
En Marzo de 2013, Phil Perviance escribe un artículo ``Don't use linksys routers'' en su blog. En
dicho artículo, Phil describe cómo le llevó sólo 30 minutos llegar a la conclusión de que
cualquier red que use un router Lynksys EA2700 es una red insegura. Uno de los ataques que
efectuó es el siguiente:

\begin{verbatim}
POST /apply.cgi
Host: 192.168.1.1
submit_button=Wireless_Basic&change_action=gozila_cgi&next_page=/etc/passwd
====>
root:x:0:0::/:/bin/sh
nobody:x:99:99:Nobody:/:/bin/nologin
sshd:x:22:22::/var/empty/sbin/nologin
admin:x:1000:1000:Admin User:/tmp/home/admin:/bin/sh
quagga:x:1001:1001:Quagga:/var/empty:/bin/nologin
firewall:x:1002:1002:Firewall:/var/empty:/bin/nologin
\end{verbatim}

\begin{enumerate}[label=\alph*)]
  \item Explicar qué tipo de vulnerabilidad queda expuesta en este ataque.
  \item ¿De qué manera podría prevenirse un ataque de este tipo?
\end{enumerate}
\end{consigna}

\paragraph{Temas.}
Clase 8 — Principios de Diseño y Vulnerabilidades (\S Catálogo de vulnerabilidades:
\emph{Injection flaws}; CWE-78 \emph{OS/Path}; \S 8 principios: mediación completa, fail-safe
defaults, mínimo privilegio).

\paragraph{Qué necesitás saber.}
El parámetro \texttt{next\_page} se envía sin validar al CGI, que lo usa como \textbf{ruta de un
archivo a leer/incluir}. El atacante pasa \texttt{next\_page=/etc/passwd} y el servidor devuelve
el contenido del archivo del sistema. Esto es una vulnerabilidad de \textbf{mezclar datos con
código / falta de control del input} --- la familia que en Clase 8 se identifica como
\emph{injection flaws} (``el 90\,\% de las vulnerabilidades''): se inyecta, a través de un input,
algo que el sistema interpreta como instrucción (acá, una ruta de archivo). El patrón concreto es
\textbf{path/directory traversal} (variante de \emph{file inclusion}), pariente cercano de
CWE-78 (OS/Command injection) en el catálogo: ``falta de control permite \dots\ un nombre de
archivo no sanitizado''. La mitigación canónica de los injection flaws en Clase 8 es
\textbf{sanitizar/validar el input} y \textbf{separar datos de código}.

\paragraph{Notas y conexiones.}
Aquí confluyen varios de los 8 principios de Saltzer \& Schroeder (Clase 8):
\begin{itemize}[leftmargin=1.4em,itemsep=1pt]
  \item \textbf{Mediación completa} (``la puerta del castillo''): \emph{toda} solicitud de un
        recurso debe pasar por un control de autorización; acá el CGI sirve un archivo arbitrario
        sin mediar control.
  \item \textbf{Fail-safe defaults / whitelist $>$ blacklist}: el default debe ser \emph{denegar};
        \texttt{next\_page} debería aceptar solo valores de una \emph{lista blanca} de páginas
        válidas, no cualquier ruta.
  \item \textbf{Mínimo privilegio}: el proceso del servidor no debería poder leer
        \texttt{/etc/passwd}; el daño se contiene si corre confinado.
\end{itemize}
Además expone \emph{sensitive data exposure} (Clase 8): revelar \texttt{/etc/passwd} (usuarios,
shells, cuentas con UID 0). Aunque las contraseñas estén en \texttt{/etc/shadow}, conocer las
cuentas habilita ataques \emph{online} sobre el login (Clase 7, \S ataques: probar root/admin).

\paragraph{Resolución.}
\textbf{(a) Vulnerabilidad.} Es un \textbf{injection flaw} por \textbf{falta de validación del
input}, concretamente \textbf{path traversal / local file inclusion}: el parámetro
\texttt{next\_page} es tomado por el CGI como una ruta del sistema de archivos sin sanitizar, lo
que permite leer cualquier archivo del dispositivo (en el ejemplo, \texttt{/etc/passwd}). Encaja
en la categoría que Clase 8 llama \emph{injection flaws} (datos tratados como código/instrucción)
y se emparenta con CWE-78. Adicionalmente constituye \emph{sensitive data exposure} (filtra la
lista de usuarios y cuentas privilegiadas del router) y viola la \textbf{mediación completa} (el
endpoint entrega un recurso sin control de autorización).

\textbf{(b) Prevención.}
\begin{enumerate}[leftmargin=1.9em,itemsep=2pt]
  \item \textbf{Validar/sanitizar el input} (mitigación canónica de injection): no usar el valor
        del cliente como ruta. Aplicar \textbf{whitelist} --- mapear \texttt{next\_page} a un
        conjunto cerrado de páginas permitidas (un índice de IDs $\to$ archivos), rechazando
        cualquier otra cosa (\emph{fail-safe defaults}). Bloquear secuencias como \texttt{../} y
        rutas absolutas.
  \item \textbf{Mediación completa}: que toda lectura de archivo pase por un control de
        autorización; el CGI no debe servir archivos fuera de su directorio web.
  \item \textbf{Mínimo privilegio / confinamiento}: ejecutar el proceso del CGI con permisos
        acotados (\emph{chroot}/jail) de modo que aunque exista el bug no alcance
        \texttt{/etc/passwd}.
  \item \textbf{No exponer información sensible} ni cuentas innecesarias en el dispositivo.
\end{enumerate}

% =====================================================================
\subsection{Ejercicio 3 — Secreto compartido de umbral $(t,n)$ y entropía}
% =====================================================================

\begin{consigna}
\begin{enumerate}[label=\arabic*)]
  \item ¿A qué se denomina \textbf{sistema de secreto compartido de umbral $(t,n)$}? Dar el
        nombre de algún sistema de secreto compartido de tipo umbral.
  \item Considerar el siguiente esquema de secreto compartido:

  \begin{quote}
  Para compartir un secreto $s \in \{0,1\}^m$ entre $n$ personas, se eligen $n-1$ valores
  $a_1, a_2, \ldots, a_{n-1}$, en forma aleatoria e independiente, donde cada
  $a_i \in \{0,1\}^m$.

  Se selecciona $a_n = s \oplus a_1 \oplus a_2 \oplus \cdots \oplus a_{n-1}$.

  Luego se distribuye $a_i$ a la $i$-ésima persona del grupo.
  \end{quote}

  Mostrar que este esquema cumple las características de un sistema de secreto compartido umbral,
  de tipo $(n,n)$ utilizando el concepto de \textbf{entropía}. (Demostrarlo para $n = 3$.)
\end{enumerate}
\end{consigna}

\paragraph{Temas.}
Clase 8 — Principios de Diseño y Vulnerabilidades (\S El secreto compartido de Shamir, en
términos de entropía; \S Flujo de información con entropía). Clase 9 — Flujo de Información y
Malware (\S Entropía de Shannon; \S Condición de flujo). Entropía como medida de información
filtrada.

\paragraph{Qué necesitás saber.}
\textbf{Entropía de Shannon} (Clase 9, \S Entropía):
\begin{equation*}
H(X) = -\sum_i p(x_i)\,\log_2 p(x_i),
\end{equation*}
mide la \textbf{incertidumbre} sobre $X$. Es \textbf{máxima} con distribución uniforme
($H=\log_2 n$) y vale $\boxed{H(X)=0 \Rightarrow \text{certeza}}$. Hay \textbf{flujo de
información} de $x$ hacia el observador cuando, al conocer algo $y$, la incertidumbre sobre $x$
\emph{baja}: $H(x\mid y) < H(x)$. \textbf{No hay flujo} (no se filtra nada del secreto) cuando la
entropía \emph{no cambia}: $H(x\mid y) = H(x)$.

\textbf{Esquema $(t,n)$ de umbral} (Clase 8, \S Shamir): hay $n$ sombras (\emph{shares}) y se
necesitan \textbf{al menos $t$} para reconstruir el secreto; con \emph{menos} de $t$ no se obtiene
\textbf{ninguna} información sobre $s$. La validez se demuestra con entropía: con $<t$ sombras
$H(S\mid\text{sombras})=H(S)$ (no hay flujo); con $\ge t$, $H\to 0$ (se recupera). El sistema
umbral por excelencia es el \textbf{esquema de Shamir} (interpolación polinómica). El esquema del
enunciado es un caso $(n,n)$ por \textbf{XOR} (one-time-pad multi-parte): hacen falta \emph{todas}
las sombras.

\paragraph{Notas y conexiones.}
Esta es la misma técnica de demostración que Clase 8 usa para Shamir, instanciada en $n=3$. La
clave es: para $s\in\{0,1\}^m$ uniforme, $H(s)=m$ bits. Como cada $a_i$ es \textbf{uniforme e
independiente}, conocer un subconjunto \emph{propio} de sombras deja a $s$ \emph{uniformemente
distribuido} (entropía intacta $=m$); solo al juntar las $n$ sombras se despeja $s$ y la entropía
cae a $0$. La propiedad del XOR ($x\oplus x = 0$, $x\oplus 0 = x$) es la que permite despejar.

\paragraph{Resolución.}

\textbf{1) Definición.} Un \textbf{sistema de secreto compartido de umbral $(t,n)$} reparte un
secreto $s$ en $n$ \emph{sombras} (\emph{shares}) entre $n$ participantes de modo que:
\begin{itemize}[leftmargin=1.4em,itemsep=1pt]
  \item con \textbf{$t$ o más} sombras se puede \textbf{reconstruir} $s$;
  \item con \textbf{menos de $t$} sombras \textbf{no se obtiene ninguna información} sobre $s$
        (la entropía de $s$ no disminuye).
\end{itemize}
Ejemplo de tipo umbral: el \textbf{esquema de Shamir} (basado en interpolación de polinomios,
$(t,n)$ general). El esquema del enunciado es un \textbf{$(n,n)$ por XOR}.

\medskip
\textbf{2) Demostración para $n=3$ (esquema $(3,3)$).} Sea $s\in\{0,1\}^m$ con distribución
\textbf{uniforme}, de modo que
\begin{equation*}
H(S) = \log_2 2^m = m \quad\text{bits}.
\end{equation*}
Se eligen $a_1, a_2$ uniformes e independientes en $\{0,1\}^m$ y se define
\begin{equation*}
a_3 = s \oplus a_1 \oplus a_2 .
\end{equation*}
Cada persona $i$ recibe $a_i$. La reconstrucción es inmediata: como el XOR es asociativo y cada
elemento es su propio inverso,
\begin{equation*}
a_1 \oplus a_2 \oplus a_3 = a_1 \oplus a_2 \oplus (s \oplus a_1 \oplus a_2) = s .
\end{equation*}

\noindent\textit{Hay que mostrar dos cosas: (i) con las 3 sombras se recupera $s$ ($H\to 0$); (ii)
con $\le 2$ sombras no se filtra nada ($H = H(S) = m$).}

\smallskip
\noindent\textbf{(i) Con las 3 sombras: $H(S\mid a_1,a_2,a_3)=0$.} Por lo anterior,
$s=a_1\oplus a_2\oplus a_3$ queda \textbf{determinado}; no hay incertidumbre:
\begin{equation*}
H(S \mid a_1, a_2, a_3) = 0 .
\end{equation*}

\smallskip
\noindent\textbf{(ii) Con 2 sombras (o menos): la entropía no cambia.} Consideremos el peor caso,
conocer dos sombras. Por simetría basta analizar un par; tomemos $\{a_1,a_2\}$ y
$\{a_1,a_3\}$:

\emph{Caso $\{a_1,a_2\}$:} $a_1,a_2$ se eligieron \emph{independientes de $s$} y de forma
uniforme. No aparecen en ninguna ecuación que ligue a $s$ (la única ecuación con $s$ es la de
$a_3$). Luego $s$ sigue siendo uniforme dado $a_1,a_2$:
\begin{equation*}
H(S \mid a_1, a_2) = H(S) = m .
\end{equation*}

\emph{Caso $\{a_1,a_3\}$ (o $\{a_2,a_3\}$):} conocemos $a_1$ y $a_3 = s\oplus a_1\oplus a_2$.
Para despejar $s = a_1\oplus a_2\oplus a_3$ \textbf{falta} $a_2$, que es uniforme e independiente y
\emph{desconocido}. Como $a_2\in\{0,1\}^m$ es uniforme, $s = (a_1\oplus a_3)\oplus a_2$ es la suma
XOR de una constante conocida con una variable uniforme $\Rightarrow$ $s$ queda \textbf{uniforme}
para el observador:
\begin{equation*}
H(S \mid a_1, a_3) = H(a_2) = m = H(S) .
\end{equation*}
(El razonamiento es idéntico al one-time-pad: XOR con un valor uniforme y secreto preserva la
distribución uniforme.)

\smallskip
\noindent\textbf{Conclusión.} Resumiendo el comportamiento de la entropía:
\begin{align*}
H(S \mid a_1)            &= H(S) = m && \text{(1 sombra: no sé nada más)}\\
H(S \mid a_1, a_2)       &= H(S) = m && \text{(2 sombras: tampoco)}\\
H(S \mid a_1, a_2, a_3)  &= 0        && \text{(3 sombras: tengo toda la información)}
\end{align*}
Con $<n=3$ sombras \textbf{no hay flujo de información} ($H$ se mantiene en $m$); con las $3$,
$H\to 0$ y se reconstruye $s$. Por lo tanto el esquema es un secreto compartido umbral de tipo
$(n,n)=(3,3)$. \qed

% =====================================================================
\subsection{Ejercicio 4 — Verdadero / Falso (biométrico, matriz vs flujo, zip bomb)}
% =====================================================================

\begin{consigna}
Indicar si las siguientes afirmaciones son verdaderas o falsas. Justificar brevemente en cada
caso.
\begin{enumerate}[label=(\arabic*)]
  \item En un sistema de autenticación biométrico, puede darse un robo de identidad.
  \item Una matriz de acceso bien diseñada evita el flujo de información no deseado.
  \item Un archivo tipo ``zip'' pequeño, por ejemplo de 42kb, que contenga 16 archivos
        comprimidos, que a su vez contenga 16 archivos comprimidos, que a su vez contiene 16
        archivos comprimidos, que a su vez contiene 16 comprimidos, que a su vez contiene 16
        archivos comprimidos, que contienen 1 archivo, con el tamaño de 4.3GB puede violar la
        seguridad del sistema al descomprimirse.
\end{enumerate}
\end{consigna}

\paragraph{Temas.}
(1) Clase 7 — Autenticación (\S Algo que soy: biométrico). (2) Clase 9 — Flujo de Información y
Malware (\S Control de acceso vs.\ flujo de información). (3) Clase 8 — Principios de Diseño y
Vulnerabilidades (\S CIA / disponibilidad: un DoS es un ataque de seguridad).

\paragraph{Qué necesitás saber.}
(1) El biométrico \textbf{no es 100\,\% eficaz}: no hay lector perfecto y un mismo patrón $C$
puede validar dos $A$ distintos (Clase 7, ``el hermano que desbloquea el celular''); además asume
que el lector no se manipula. (2) El control de acceso (la matriz) limita \emph{operaciones sobre
objetos}, no el \emph{destino} de la información una vez leída; un sujeto con permiso de lectura
puede copiar el dato a otro objeto legible por terceros --- la matriz \textbf{no} controla el
flujo (Clase 9). (3) El primer requisito de seguridad es la \textbf{disponibilidad}: un
\textbf{DoS es un ataque de seguridad} (Clase 8). Una \emph{zip bomb} agota CPU/disco/memoria al
descomprimir.

\paragraph{Notas y conexiones.}
(2) es un error clásico que el índice marca explícitamente como penalizable: ``afirmar que el
control de acceso/matriz \emph{evita} el flujo''. La matriz es un \textbf{mecanismo abierto}
(Clase 9): asegura una parte, no todo; debe complementarse con control de \emph{flujo}
(compile-time o run-time) y aislamiento. (3) conecta con la tríada CIA y con malware/payload: la
zip bomb es un payload de \emph{disponibilidad}.

\paragraph{Resolución.}
\begin{enumerate}[label=(\arabic*),leftmargin=2.2em,itemsep=4pt]
  \item \textbf{VERDADERO.} La biometría no es perfecta: no existe lector ideal y un mismo
        registro almacenado $C$ puede aceptar información $A$ de otra persona (falsos positivos;
        ej.\ el familiar que desbloquea el teléfono), y se baja el umbral de error para reducir
        falsos rechazos, lo que \emph{aumenta} el riesgo de aceptar a otro. Además se asume que el
        lector no se puede manipular/suplantar (foto ante un lector de cara, intercepción del
        sensor). Cualquiera de estas vías permite que un atacante sea autenticado como otra
        identidad: \textbf{sí puede darse un robo de identidad}.
  \item \textbf{FALSO.} Una matriz de acceso (control de acceso) controla \emph{qué operaciones}
        puede hacer cada sujeto sobre cada objeto, pero \textbf{no controla el flujo} de la
        información una vez leída. Ejemplo (Clase 9): un empleado con permiso de \emph{leer} su
        sueldo puede \emph{escribirlo} en un archivo de notas que otros leen --- cada sentencia
        respeta la matriz, pero la información \emph{fluye} a un tercero. El control de acceso es
        un \emph{mecanismo abierto}: limita el acceso a objetos, no el destino de los datos. Para
        evitar el flujo no deseado se necesita \textbf{control de flujo} (compile-time o run-time)
        y/o aislamiento, no solo la matriz.
  \item \textbf{VERDADERO.} Es una \textbf{zip bomb} (bomba de descompresión): un archivo
        diminuto que, por compresión anidada, se expande a un tamaño enorme. Al descomprimirse
        recursivamente agota recursos del sistema (CPU, memoria, disco), provocando una
        \textbf{denegación de servicio (DoS)}. Como el primer requisito de seguridad es la
        \emph{disponibilidad}, un DoS \textbf{es} un ataque a la seguridad del sistema: el archivo
        \emph{sí} puede violar la seguridad al descomprimirse. La mitigación es validar/limitar el
        ratio y la profundidad de descompresión (mediación completa sobre el recurso).
\end{enumerate}

% =====================================================================
\subsection{Ejercicio 5 — Módulos de Control de Acceso y Autenticación}
% =====================================================================

\begin{consigna}
\textit{[El gráfico muestra: }sujeto $\to$ \textbf{Módulo de \ldots} $\to$ \textbf{Módulo de
\ldots} $\to$ objeto.\textit{]}

El gráfico representa en forma simplificada los módulos que regulan el acceso de un usuario a un
objeto protegido.
\begin{enumerate}[label=\alph*)]
  \item Completar el gráfico para indicar cuál sería el Módulo de Control de Acceso y cuál el de
        Autenticación.
  \item Mencionar una vulnerabilidad que deba tenerse en cuenta en el Módulo de Control de Acceso.
        Explicar.
  \item Mencionar una vulnerabilidad que deba tenerse en cuenta en el Módulo de Autenticación.
        Explicar.
\end{enumerate}
\end{consigna}

\paragraph{Temas.}
Clase 7 — Autenticación (\S Autenticación $\neq$ Autorización; \S Sistema de autenticación
$\langle A,C,F,L,S\rangle$; \S Ataques). Clase 6 — Políticas y Control de Acceso (\S DAC: ACL,
conflictos de reglas). Orden de los módulos: primero se autentica (¿quién sos?), luego se
autoriza el acceso (¿podés acceder a \emph{este} objeto?).

\paragraph{Qué necesitás saber.}
\textbf{Autenticación $\neq$ autorización} (Clase 7): autenticación responde ``¿quién sos?''
(verifica identidad); autorización/control de acceso responde ``ya autenticado, ¿podés acceder a
\emph{este} recurso?''. El orden lógico es: el sujeto se \textbf{autentica} primero, y recién con
una identidad establecida se aplica el \textbf{control de acceso} sobre el objeto. Por eso, en el
flujo sujeto $\to$ \fbox{1} $\to$ \fbox{2} $\to$ objeto, el módulo \fbox{1} (pegado al sujeto) es
\textbf{Autenticación} y el \fbox{2} (pegado al objeto) es \textbf{Control de Acceso}.

\paragraph{Notas y conexiones.}
El sistema de autenticación se modela como la tupla $\langle A,C,F,L,S\rangle$ (Clase 7): $A$ info
del usuario, $C$ info complementaria que guarda el sistema (p.\ ej.\ el hash), $F$ \emph{deriva}
($A\to C$, no invertible), $L$ \emph{corrobora} ($A\times C\to\{0,1\}$, el login), $S$ funciones
de alta/cambio. Los ataques \emph{offline} atacan $F$ (robaron $C$ y prueban $a$ hasta $F(a)=c$);
los \emph{online} atacan $L$ (login con cuentas conocidas). El control de acceso, en cambio, falla
por reglas mal resueltas (ACL con conflictos, Clase 6) o por no mediar todos los accesos.

\paragraph{Resolución.}

\textbf{a) Completar el gráfico.} El sujeto primero debe probar su identidad y luego pedir acceso
al objeto:
\begin{center}
\begin{tikzpicture}[font=\small,>=Stealth,node distance=10mm]
  \node[circle,draw,fill=itbaBlue!30,minimum size=0.9cm] (s) {sujeto};
  \node[draw,rounded corners,fill=boxBlue,minimum width=2.6cm,minimum height=1.1cm,right=of s,align=center]
       (auth) {Módulo de\\\textbf{Autenticación}};
  \node[draw,rounded corners,fill=boxBlue,minimum width=2.6cm,minimum height=1.1cm,right=of auth,align=center]
       (ac) {Módulo de\\\textbf{Control de Acceso}};
  \node[circle,draw,fill=slideGray,minimum size=0.9cm,right=of ac] (o) {objeto};
  \draw[->,line width=1pt] (s) -- (auth);
  \draw[->,line width=1pt] (auth) -- (ac);
  \draw[->,line width=1pt] (ac) -- (o);
\end{tikzpicture}
\end{center}
\textbf{Justificación:} primero se \emph{autentica} (``¿quién sos?''), estableciendo la identidad;
recién entonces el \emph{control de acceso} decide si esa identidad está autorizada a operar sobre
\emph{ese} objeto (``ya autenticado, ¿podés acceder a este recurso?''). Autenticación $\neq$
autorización.

\textbf{b) Vulnerabilidad del Módulo de Control de Acceso.} \emph{Resolución incorrecta de reglas
en una ACL con conflictos} (Clase 6). Si una política tiene reglas que se contradicen (una permite,
otra deniega) y no se define bien la estrategia de resolución (unión / Grant-All = intersección /
first-match / best-match / deny-over-allow), un sujeto puede terminar con \textbf{más permisos de
los previstos}. Otra falla central es no garantizar \textbf{mediación completa}: si algún camino
de acceso al objeto \emph{no} pasa por el módulo, el control se puede saltear. En ambos casos el
resultado es una \textbf{elevación de privilegios / acceso no autorizado} al objeto protegido.

\textbf{c) Vulnerabilidad del Módulo de Autenticación.} \emph{Ataque a las credenciales} (Clase 7).
Si el sistema guarda mal $C$ (texto plano, sin salting, hash débil) un atacante que obtiene la base
hace un \textbf{ataque offline} contra $F$: prueba candidatos $a$ hasta hallar $F(a)=c$ (diccionario,
fuerza bruta, rainbow tables). Si el módulo de login $L$ no se protege, sufre \textbf{ataques
online} (fuerza bruta / credential stuffing sobre cuentas conocidas como root/admin) por falta de
\emph{exponential back-off} o bloqueo. El efecto es la \textbf{suplantación de identidad}: el
atacante se autentica como otro usuario y desde ahí el control de acceso lo trata como legítimo.
(Conexión con Ej.\ 4.1: en biometría, la vulnerabilidad análoga es el falso positivo / lector
manipulable.)
